\documentclass[10pt, aps, rmp,  twocolumn, superscriptaddress, floatfix, showkeys]{revtex4-2}
\usepackage[shortlabels]{enumitem}
\bibliographystyle{plainnat}
%\usepackage[latin1]{inputenc}
\usepackage{multirow}
\usepackage{booktabs}
\usepackage{amsmath}
\usepackage{amsfonts}
\usepackage{graphicx}
\usepackage{calrsfs} 
\usepackage{times}
\usepackage[sort&compress]{natbib}
\usepackage{subfigure}
\usepackage{xcolor}
\usepackage{calrsfs}
\usepackage{soul}
\usepackage{booktabs, multirow}
\usepackage{comment}
\usepackage{verbatim}
\usepackage{framed}
\usepackage{lineno}
\linenumbers

\usepackage{xcolor}
\usepackage{mdframed}

\newmdenv[
  backgroundcolor=gray!10,
  linecolor=gray!10,
  linewidth=0pt,
  innerleftmargin=10pt,
  innerrightmargin=10pt,
  innertopmargin=10pt,
  innerbottommargin=10pt,
  skipabove=10pt,
  skipbelow=10pt
]{graybox}

% Preamble code to extract detailed section word counts
\newcommand{\detailtexcount}[1]{%
  \immediate\write18{texcount -v -merge -q #1.tex > #1.wcdet}%
  \verbatiminput{#1.wcdet}%
}

\providecommand{\keywords}[1]{\textbf{\textit{Keywords---}} #1}


\begin{document}
\newcommand{\fz}[1]{\textcolor{blue}{#1}}
\newcommand{\tbd}[1]{\textcolor{blue}{[TBD] #1}}

%\title{Alternative genome instability states in a digital model of tumorigenesis with whole genome duplication}

%\title{Evolution of critical cancer instability after whole genome duplication}
%\title{Evolutionary paths towards critical cancer instability}
%\title{Self-Organized cancer fragility emerges from chromosomal instability}
\title{A Probabilistic Ensemble-Based Bayesian Method for Mechanistic GRN Inference}
%\title{Genomically unstable evolution drives tumours to the edge of viability}

  
\providecommand{\DFA}{Department of Physics and Astronomy ``Galileo Galilei'', University of Padova, Padova, Italy}
\providecommand{\PNC}{Padova Neuroscience Center, University of Padua, Via G. Orus 2, 35131 Padova, Italy}
\providecommand{\PCNM}{Padua Center for Network Medicine, University of Padua, Via F. Marzolo 8, 315126 Padova, Italy}
\providecommand{\INFN}{Istituto Nazionale di Fisica Nucleare, Sez. Padova, Italy}

\author{Francesco Zambelli}
\affiliation{\DFA}
\author{Valeria D'Andrea}
\affiliation{\DFA}
\author{Manlio De Domenico}
\affiliation{\DFA}
\affiliation{\PNC}
\affiliation{\PCNM}
\affiliation{\INFN}

\vspace{0.4 cm}
\begin{abstract}
\vspace{0.2 cm}  
Recovering the interaction structure of a complex stochastic dynamical system from time-series observations is an inherently ill-posed problem: the posterior over network configurations is high-dimensional, multimodal, and shaped by the non-equilibrium nature of the available data. We propose a Bayesian ensemble inference framework grounded in the Onsager-Machlup path probability formalism, which provides a rigorous likelihood for SDE-governed trajectories without requiring stationarity, and is coupled with a node-wise block-coordinate Metropolis-Hastings sampler that exploits the factorization of the path probability to make inference tractable. A evidence-estimation-based model comparison scheme jointly recovers the dynamical parameters alongside the network topology. Applying the framework to gene regulatory network (GRN) inference, we demonstrate on synthetic data that ensemble-based posterior analysis produce more reliable results, in particular when considering perturbation predictions and network topological statistics, than single-point estimation. The framework is general and can be applied to any mechanistic SDE model, offering a powerful tool for uncertainty-aware inference in complex stochastic systems.
\vspace{0.2 cm}
\end{abstract}

\keywords{Cancer, Tumor Growth, Genomic Instability, Error Threshold}

\maketitle

%Paragraph 1: The General Problem. Reconstructing the causal topology of non-linear networked systems from sparse time-series data is a universal challenge across disciplines (e.g., Gene Regulatory Networks, epidemiological models, ecological population dynamics). Introduce these systems as being characterized by two things: saturating behaviors (carrying capacity/resource limits) and intrinsic stochasticity.

%Paragraph 2: The Data Limitations. In reality, time-series data from these systems is sparse, noisy, and often measured at large intervals.

%Paragraph 3: Current Approaches vs. SDEs. Explain why traditional methods fail: Correlation methods (like WGCNA/Time-Delayed Pearson) only capture symmetric linear relationships and cannot distinguish direct causality from indirect cascades. Deterministic ODEs overfit because they assume all variance is structural; when fit to noisy sparse data, the optimizer hallucinates edges to explain random fluctuations. Introduce Stochastic Differential Equations (SDEs) as the rigorous solution because they explicitly separate deterministic structural dynamics (the network) from intrinsic stochastic noise.

%Paragraph 4: Paper Roadmap. Outline the transition from MAP estimates to Bayesian ensembles, and the mathematical innovations required to make it computationally feasible.


\section{Introduction}

\section{Bayesian Network Inference}

%Paragraph 1: The Inverse Problem. Formulate network reconstruction as an inverse problem. Introduce Bayes' Theorem.




When trying to approach the inverse problem of network reconstruction from some data $\mathbf{X}$, the most rigorous and transparent approach is using Bayesian inference:

\begin{equation}
    P(\mathbf{W},\theta|\mathbf{X}) = \frac{P(\mathbf{X}|\mathbf{W}, \theta)P(\mathbf{W}, \theta)}{P(\mathbf{X})}
\end{equation}


%Paragraph 2: Components of Bayes. Explain the Likelihood (how well the network explains the trajectory) and the Prior (our biological assumptions, e.g., sparsity).
Using this equation, we suppose that the data we have are generated by a specific generative model, described by the likelihood function $P(\mathbf{X}|\mathbf{W}, \theta)$, where $\mathbf{W}$ is the encodes the network structure and weights and $\theta$ are the extra parameters of the generative model. 
In this work we will focus on data organized in time series, so, if we have a network of $N$ nodes, and the data are organized in $M$ time points, we can write the data as a matrix $\mathbf{X} = \{\mathbf{x}_{t_1}, \mathbf{x}_{t_2}, ..., \mathbf{x}_{t_M}\}$ of size $M \times N$, where $\mathbf{x}_m$ is a vector of size $N$ that encodes the state of the nodes at time $t_m$. In this case, the likelihood function can be factorized as:

\begin{equation}
    P(\mathbf{X}|\mathbf{W}, \theta) = \prod_{m=1}^{M-1} P(\mathbf{x}_{t_m+\Delta t}|\mathbf{x}_{t_m}, \mathbf{W}, \theta)
    \label{eq:general_likelihood}
\end{equation}

In most of the realistic scenarios, the data we have are scarce and noisy, and the generative models are complex and high-dimensional. For this reasons, when we try to perform inference on the generative model, we can encounter two main problems: non-identifiability and overfitting. Non-identifiability occurs when multiple distinct configurations of the interactions can produce trajectories that are statistically indistinguishable from the data. Overfitting occurs when the model is too complex and fits the noise in the data rather than the underlying signal. 

To overcome these issues we need to incorporate an extra term which acts as a regularizer, penalizing complex models and incorporating prior knowledge about the system. This term is the prior distribution $P(\mathbf{W}, \theta)$, that enables the inference process to be guided towards more plausible configurations, and to avoid overfitting.


%Paragraph 3: MDL and Compression. Frame Bayesian optimization conceptually: finding the network that maximizes Minimum Description Length (MDL) compression of the data.

One way to interpret the Bayesian inference framework is through the lens of Minimum Description Length (MDL) principle. Taking the logarithm of Bayes' theorem, we can rewrite it as:

\begin{equation}
    \log P(\mathbf{W},\theta|\mathbf{X}) = \log P(\mathbf{X}|\mathbf{W}, \theta) + \log P(\mathbf{W}, \theta) - \log P(\mathbf{X})
\end{equation}

In this way the inference problem can be seen as a maximum compression problem, where we want to find the network configuration that requires the least amount of information to describe the data. This is a crucial aspect of Bayesian inference, as it clearly states the goal of the inference process: to find the network configuration that best explains the data while being as simple as possible.

%Paragraph 4: The Paradigm Shift. Contrast finding the single "best" network (MAP estimate) with generating a probabilistic ensemble. Why ensembles are necessary when single false-positive edges can ruin biological interpretations.

Once we have defined both the generative model and the prior distribution, we can find the configuration that is able to best compress the data, which is called the Maximum A Posteriori (MAP) estimate. However, in many cases, the MAP estimate is not sufficient to capture the full complexity of the inference problem. In particular, the best network configuration often does not correspond exactly to the true underlying network, due to data scarcity, noise and other aspects that cannot be completely comptured by the generative model. 

This aspect is particularly important in context where the mis-specification of a single edge can have a relevant impact on the subsequent system's analysis. This is particulary relevant in the context of gene regulatory networks, where in most of the case, the inferred network is used to perform perturbation analysis, and the mis-specification of a single edge can lead to completely wrong predictions. 

For this reason, it is crucial to consider the full posterior distribution over network configurations, rather than just the MAP estimate. This allows us to capture the uncertainty in the inference process and to generate a probabilistic ensemble of plausible network configurations, which can be used to perform more robust and reliable analyses. This can be acheived by sampling from the posterior distribution using Markov Chain Monte Carlo (MCMC) methods, which allow us to generate a representative set of network configurations that are consistent with the generative model and the model's constraints.


\subsection{Capturing Saturation and Relaxation in Networked Systems}


%Paragraph 1: The SDE Model. Present the specific SDE: continuous exponential decay, saturating non-linear activation (sigmoid), and the interaction matrix W. Crucially, explain why we chose this model: It is the "simplest non-trivial model" capable of capturing the universal features of bounded complex networks—namely, intrinsic exponential relaxation (without which states would blow up) and bounded non-linear coupling (because real-world networks exert saturating, not infinite, influence).

As a generative model for the data, we consider a system of stochastic differential equations (SDEs) that describe the dynamics of a network of interacting nodes. This is the most general and flexible framework for modeling complex systems, as it allows us to capture both the deterministic dynamics of the system and the intrinsic stochasticity that characterizes the aquired data. 
As a representative case study, we focus on the following model:

\begin{equation}
    dx_i = -\gamma \, x_i + V \sigma \, \Bigg(\sum_{j=1}^{N} W_{ij}\, x_j \Bigg) +  \sqrt{D \, x_i} \, dW_i(t)
\end{equation}

Altough simple, this model is able to capture the universal features of bounded complex networks, namely, intrinsic exponential relaxation and bounded non-linear coupling. Such features are common to many real-world systems, including brain networks, ecological systems, macroeconomics and many others. A particular example of such systems is represented by gene regulatory networks, where the nodes represent genes and the edges represent regulatory interactions between them. In this context, the exponential decay term represents the natural degradation of gene products, while the saturating non-linear activation term represents the limited capacity of genes to influence each other. The interaction matrix $W$ encodes the strength and directionality of these interactions, and the multiplicative noise term captures the intrinsic stochasticity that characterizes gene expression. For the following, we will use as a interpretative framework the context of gene regulatory networks.

In this kind of systems, the correct inference of the interaction matrix $W$ is crucial for understanding the underlying regulatory mechanisms and for making predictions about the system's behavior under different conditions, and for this reason, the Bayesian inference framework is particularly well-suited for this task, as it allows to produce a probabilistic ensemble of plausible network configurations that can be used to perform more robust and reliable analyses.


%Paragraph 2: Multiplicative Noise. Argue the choice of Cox-Ingersoll-Ross (CIR) multiplicative noise over additive noise, tying it to realistic biological burstiness.

A particular aspect of the model that we consider is the choice of the multiplicative noise term, which is modeled as a Cox-Ingersoll-Ross (CIR) process. This choice is motivated by the fact that gene expression is known to be a bursty process, where the production of gene products occurs in short bursts rather than at a constant rate. This burstiness can be captured by the multiplicative noise term, which allows for fluctuations in the gene expression levels that are proportional to the current state of the system. In contrast, additive noise would not capture this feature, as it would introduce fluctuations that are independent of the current state of the system.

%Paragraph 3: Limitations of Saturation. Discuss the difficulty of inferring saturating dynamics from sparse time-series (if a gene is saturated, changes in input are invisible).




%Paragraph 4: Ground Truth Generation. Explain the choice of a scale-free topology for generating the synthetic validation data and parameter selection.

\subsection{Beyond Euler-Maruyama: Exact Transition Moments for Large Measurement Intervals}

%Paragraph 1: The Likelihood Bottleneck. Explain that Bayesian inference requires evaluating the probability of discrete trajectory points. Introduce the standard Euler-Maruyama (E-M) discretization.

To compute the total likelihood of a trajectory (Eq.~\ref{eq:general_likelihood}), the stohastic process is treated as a Markov chain, requiring the evaluation of the transition probability $P(\mathbf{x}_{t_m+\Delta t}|\mathbf{x}_{t_m}, \mathbf{W}, \theta)$ for every step.


Traditionally, this transition probability is approximated using the Euler-Maruyama numerical scheme. This approach assumes that the entire drift derivative remains constant over the measurement interval $\Delta t$, yielding a Gaussian transition density:

\begin{equation}
    P(x_{t+\Delta t} \mid \mathbf{X}_t) \sim \mathbf{N}\left( x_t + \Delta t \big(-\gamma x_t + V\sigma(m_t)\big), \; D x_t \Delta t \right)
\end{equation}

where $m_t = \sum_{j=1}^{N} W_{ij}\, x_j(t)$ is the aggregate regulatory input to node $i$ at time $t$.


%Paragraph 2: The Discretization Penalty. Detail why E-M fails for biological data: large sampling intervals cause massive tangent overshoots, forcing the optimizer to break the network topology to compensate.

While computationally trivial, this approximation introduces catastrophic structural bias when applied to sparse biological datasets. Real-world transcriptomic time-series often exhibit large measurement intervals (e.g., hours between sampling). Under large $\Delta t$, the Euler-Maruyama linear tangent completely ignores the exponential relaxation inherent to the $-\gamma x_t$ degradation term. This results in severe mathematical overshoot, where the predicted mean may drop below zero or artificially inflate.

%Paragraph 3: The Exact Transition Solution. Introduce the Zero-Order Hold assumption on the regulatory input. Show how this isolates a linear SDE that can be solved exactly.

To mitigate the discretization error without incurring the computationally intractable cost of inferring unobserved latent variables between measurement points, we introduce an exact analytical integration approach predicated on a zero-order hold assumption.

We assume that the aggregate regulatory input from the rest of the network, $m_i(t)$, remains piecewise constant over the short interval $[t, t+\Delta t]$. Under this assumption, the non-linear coupling term $V \, \sigma(m_i(t))$ simplifies to a constant scalar $b$.

The governing SDE for a single gene isolated during this time step reduces to a strictly linear stochastic differential equation:

\[ dx_t = (-\gamma x_t + b) dt + \sqrt{D x_t} dW_t \]

By isolating the gene's state variable $x_t$, the SDE is mathematically equivalent to a continuous CIR process, which permits an exact analytical solution for its transition moments regardless of the magnitude of $\Delta t$.

%Paragraph 4: Theoretical Bounds. Briefly report the exact mean (exponential relaxation) and exact variance (via Itô's Isometry), proving it eliminates the discretization boundary constraints.

Here is the concise, standalone mathematical block you can drop directly into the main text of your article. It gives the reader the exact formulas they need to implement the likelihood, while leaving the heavy stochastic calculus (Itô's Isometry) for the Supplementary Material.

The exact transition probability density for the state of node $i$ at time $t+\Delta t$ can be approximated by a Gaussian $\mathcal{N}(\mu, \Sigma)$, parameterized by its exact analytical moments. 
    
By defining the deterministic steady-state asymptote as $\mu_\infty = \frac{V \, \sigma(m_i(t))}{\gamma}$, the exact expected mean is given by the exponential relaxation:

\[ \mu(t+\Delta t) = \mu_\infty + \big(x_i(t) - \mu_\infty\big)e^{-\gamma_i \Delta t} \]

And the exact integrated variance accumulated over the interval $\Delta t$ is given by:

\begin{equation}
    \begin{aligned}
    \Sigma(t+\Delta t) &= \frac{D_i}{\gamma_i}\big(x_i(t) - \mu_\infty\big)\Big(e^{-\gamma_i \Delta t} - e^{-2\gamma_i \Delta t}\Big) \\
    &\quad + \frac{D_i \mu_\infty}{2\gamma_i}\Big(1 - e^{-2\gamma_i \Delta t}\Big)
    \end{aligned}
\end{equation}

The full analytical derivation of these moments via Itô's Isometry is provided in the Supplementary Material.

\subsection{Inference engine}

%Paragraph 1: Graph-Tool Framework. Present Tiago Peixoto’s perspective on network inference (stochastic block models, non-parametric priors).

Once the generative model is defined, we employ an extension of the Markov Chain Monte Carlo (MCMC) inference engine developed by Peixoto~\cite{peixoto2025}, implemented via the \textit{graph-tool} library. 
This method incorporates a hierarchical, non-parametric Minimum Description Length (MDL) prior, which goal is to rigorously penalize model complexity without relying on weight shrinkage. It achieves this via \textbf{weight quantization}: rather than treating every non-zero edge as having an infinitely precise, unique weight (which would require infinite information to describe), the algorithm is forced to compress the network's continuous weights into a highly redundant, finite vocabulary of discrete categories. 

This adaptive compression is formalized by marginalizing the prior over several latent structural and parametric variables:
\begin{equation}
    \begin{aligned}
        P(\mathbf{W} \mid \lambda, \Delta) &= \sum_{z,K,\mathbf{A},E} P(\mathbf{W} \mid z, \mathbf{A})\, P(z \mid \lambda, \Delta, K)\, \\
        &\quad\times P(K \mid \mathbf{A})\,P(\mathbf{A} \mid E)\,P(E)
    \end{aligned}
    \label{eq:mdl_prior}
\end{equation}

where $\mathbf{W}$ represents the continuous interaction weights of the network, and $\mathbf{A}$ is its corresponding binary adjacency matrix containing exactly $E$ non-zero edges. The parameter $K$ denotes the total number of unique, discrete weight categories into which the edges are grouped, while $z = \{z_1, \dots, z_K\}$ represents the set of continuous values for those $K$ categories.

This hierarchical formulation allows the inference engine to dynamically adjust the complexity of the model. Several of the marginal distributions represent simple uniform priors that serve as baseline complexity penalties: $P(E)$ is uniformly distributed over the possible number of edges and $P(K \mid \mathbf{A})$ is uniform over the number of discrete weight categories ($K \le E$).

The term responsible for the sparse structure of the network is:

\begin{equation} 
    P(\mathbf{A} \mid E) = \frac{\delta_{\sum_{i,j}A_{ij}, E}}{\binom{N^2}{E}} 
\end{equation}

which penalizes topologies with large number of edges, leading to small description length for sparse graphs.

The other two terms govern how the continuous weights are compressed into the discrete vocabulary $z$. The first one, $P(\mathbf{W} \mid z, \mathbf{A})$, quantifies how much the number of paramters categories impact on the compression level of the data.
If $m_k$ represents the number of edges assigned to a specific weight category $k$, this probability distribution yields a multinomial combinatorial term proportional to $\frac{\prod_k m_k!}{E!}$. From an MDL perspective, this term heavily rewards parameter reuse. Because the inverse of the multinomial coefficient maximizes when many edges share the exact same weight category (e.g., $m_k$ is large), the algorithm is strongly incentivized to minimize $K$ and group the weights into a small, highly redundant vocabulary.

The second term, $P(z \mid \lambda, \Delta, K)$, prevents the absolute values of the weight categories from diverging. They are drawn from a quantized Laplace distribution defined by the scale parameter $\lambda$ and a numerical precision $\Delta$, yielding a probability proportional to $e^{-\lambda \sum_{k=1}^K |z_k|}$. This does not yeald to strong shrinkage of the weights, making it a more robust alternative to the common LASSO regularization.

Alongside to these terms, the prior could also include the Stochastic Block Model (SBM) with degree correction, which allows the algorithm to exploit the modular structure of the network to further compress the description length. This is particularly useful in the context of gene regulatory networks, where genes are often organized into functional modules that share similar regulatory patterns, and for the case scenario considered in this work, is able to classify between the master regulators and the downstream targets, improving the overall inference performance.

All these elements compose a non-parametric prior that is able to adaptively penalize model complexity, while still allowing for a rich and flexible representation of the network structure and weights. This is crucial for avoiding overfitting and ensuring that the inferred networks are both biologically plausible and statistically robust.

%Paragraph 2: Integration. Explain how your exact analytical likelihood is embedded into the MCMC engine of graph-tool.

To optimize the joint posterior of the network topology and the global kinetic parameters, we employ a block-coordinate optimization scheme. The algorithm alternates between two phases: first, the discrete topological space (edges and block assignments) is sampled stochastically using Metropolis-Hastings MCMC sweeps. The sampling scheme for the topology is based on the $NN_{Descent}$ algorithm, which is able to solve the combinatorial optimization problem in a subquadratic time complexity of order $O(N\,log^2 N)$, instead of the naive $O(N^2)$ time complexity for standard gradient descent algorithms. On the other hand, the classes of parameters values are obtained by a system of merging and splitting of the weight classes and random bisection to obtain new values for the weight categories.

Second, the global continuous kinetic parameters $\theta = \{\gamma, V, D\}$ are deterministically optimized conditional on the current topology using the Nelder-Mead simplex algorithm to minimize the total description length. This choice was made because the likelihood function is convex with respect to the kinetic parameters, and given a fixed topology and weight categories, the optimization problem is well-posed and can be solved efficiently using a deterministic algorithm.

This hybrid approach allows the engine to efficiently navigate the highly combinatorial structural space while quickly relaxing the continuous dynamics into their local optimum.

\section{Results}

\subsection{The necessity of simple models: the problem of overfitting}

In the SDE model we consider (Eq.~1), the dynamical parameters $\theta = \{\gamma, V, D\}$ are global parameters that are shared across all nodes in the network. This means that the inference of these parameters is not node-specific, but rather a global optimization problem that takes into account the entire network's dynamics. 

That is a fantastic observation, and it is a fundamental problem in inverse modeling known as **parameter degeneracy** (or non-identifiability). What you experienced is exactly what happens when a model is overparameterized relative to the sparsity of the data.

A fundamental challenge in inferring SDE-driven networks from sparse time-series data is the degeneracy between intrinsic node dynamics and network-driven coupling. If the kinetic parameters ($\gamma, V, D$) are allowed to vary freely for every individual node, the model becomes overparameterized. Under the MDL framework, the optimizer naturally penalizes structural complexity; thus, it preferentially tunes the local intrinsic parameters to absorb the trajectory variance, resulting in an artificial collapse of the network topology (an empty graph).

To break this degeneracy, we constrain the kinetic parameters to be global, shared symmetrically across all nodes in the network. By stripping the nodes of their individual parametric flexibility, we force the intrinsic dynamics to act as a uniform, homogeneous baseline. Consequently, the only mathematical avenue left for the inference engine to explain the diverse, node-specific heterogeneous fluctuations observed in the data is to reconstruct the true underlying interaction topology $W$. This global constraint acts as a powerful structural inducer, guaranteeing that the inferred network is driven by genuine interactions rather than localized parametric overfitting.

In Fig.~\ref{fig:degeneracy}, we illustrate this phenomenon. When the kinetic parameters are allowed to vary per node, the inferred network collapses to a trivial empty graph, as the optimizer absorbs all variance into the local parameters. In contrast, when the kinetic parameters are constrained globally, the inference engine successfully recovers a rich and biologically plausible network topology.
% maybe it will depends on the sample size

It is possible to show that this kind of model are able to help to avoid overfitting also with respect to other inference models. We show in Fig.~\ref{fig:comparison} a comparison between the inference performance of our model and other inference models, ...

\subsection{Stochasticity impose strict discretizations boundaries}

% Present the "Stress Discretization" benchmarking results. Show how the naive E-M approach collapses as \delta t increases, while your Analytical Transition approach maintains robust AUPRC/AUROC scores.

As discussed in the previous sections, it is necessary to overcome the discretization error introduced by the Euler-Maruyama approximation, to deal with the large measurement intervals that are typical of biological data. In Fig.~\ref{fig:discretization}, we show the impact of the discretization error on the inference performance, by comparing the inference results obtained using the Euler-Maruyama approximation and the exact analytical integration approach.


\subsection{Working outside stationarity conveys more information}

% Demonstrate that inference relies on transient dynamics. If the GRN is stuck at a steady-state, highly correlated genes are indistinguishable from causal genes. Show that trajectories actively relaxing to steady-state provide the necessary causal information.

\subsection{Toward an ensemble perspective}

\paragraph{Probabilistic evaluation of inference results}

% Show that the MCMC produces a distribution of networks. Define edge probabilities rather than binary present/absent edges.

\paragraph{Impact on perturbation analysis}

% The ultimate biological payoff: demonstrate that predicting the effect of a drug/knockout (perturbation) is highly sensitive to single edges. Show how an ensemble of networks gives a confidence interval for a perturbation outcome, whereas a single MAP network gives a brittle, potentially wrong prediction.

\section{Conclusions}

\vspace{0.2 cm}

{\bf Acknowledgments}
F.Z. is supported by a Decreto Ministeriale n. 118 del 02/03/2023, M4C1 I. 4.1 from Piano Nazionale di Ripresa e Resilienza (PNRR). V.d.A. and M.D.D. acknowledge financial support from the EU funding within the MUR PNRR “National Center for HPC, BIG DATA AND QUANTUM COMPUTING” (Project no. CN00000013 CN1) and from the INFN grant “LINCOLN”. The authors acknowledge the generous support of CloudVeneto for the computing infrastructure.

\vspace{0.2 cm}
\vspace{0.2 cm}


{\bf Author contributions} 



\vspace{0.2 cm}



\renewcommand{\bibsection}{\section{References}}
\bibliographystyle{plainnat}
\bibliography{biblio}

\end{document}